% Multi-scale data assimilation for the LES ensembles: a design note.
% Build: latexmk -pdf -outdir=<tmp> docs/research/multiscale_nested_da.tex
\documentclass[11pt,a4paper]{article}
\usepackage[margin=2.5cm]{geometry}
\usepackage{amsmath,amssymb,bm}
\usepackage{booktabs}
\usepackage[hidelinks]{hyperref}

\newcommand{\T}{^{\top}}
\newcommand{\R}{\mathbb{R}}

\title{Multi-scale data assimilation for urban LES:\\
parameters, inflow and velocity field}
\author{Design note, pyurbanair}
\date{November 2026}

\begin{document}
\maketitle

\begin{abstract}
A three-time-scale design for estimating closure and roughness parameters, the
time-varying inflow and the velocity field from sensor data, inspired by the
nested filters of P\'erez-Vieites, Molina-Bulla \& M\'iguez (arXiv:2204.07795).
It keeps their factorisation of the posterior. It does not adopt their
algorithms, which would cost $N \times J$ LES runs per step and need a UKF or
EKF on the turbulent field. Instead there are two stages: a joint ESMDA on
parameters and inflow against different statistics of the observations,
followed by a streamed, localised ensemble smoother on the saved velocity
fields. Stage 2's local and global updates have had one small twin test
(Section~\ref{sec:twin}); the rest are untested proposals.
\end{abstract}

\section{Variables and time scales}

\begin{center}
\begin{tabular}{@{}llll@{}}
\toprule
 & Variables & Time scale & Paper's role \\
\midrule
$\bm\theta$ & $b_0, b_1, b_2, z_0$ & constant & static parameters \\
$\bm x$ & inflow angle and speed (AR2 knots) & $\sim 60$--$200$\,s & slow state \\
$\bm z$ & velocity field ($9.6\times10^4$ values, \S\ref{sec:sizes}) & $\sim 1$\,s & fast state \\
\bottomrule
\end{tabular}
\end{center}

\noindent The paper factorises the posterior over one window as
\begin{equation}
  p(\bm z, \bm x, \bm\theta \mid \bm y)
  = p(\bm z \mid \bm x, \bm\theta, \bm y)\,
    p(\bm x \mid \bm\theta, \bm y)\,
    p(\bm\theta \mid \bm y),
  \label{eq:factor}
\end{equation}
and obtains the likelihood of the slower variables by integrating out the
faster ones:
\begin{equation}
  p(\bm y \mid \bm x, \bm\theta)
  = \int p(\bm y \mid \bm z, \bm x, \bm\theta)\,
         p(\bm z \mid \bm x, \bm\theta)\, d\bm z .
  \label{eq:marginal}
\end{equation}

\subsection*{Where the LES differs from the paper's model}
\begin{enumerate}
\item \textbf{One-way coupling.} In the paper the slow state is driven by the
  time average of the fast state, $\bm x_t = \bm x_{t-1} + \Delta_x(f_x +
  g_x(\bar{\bm z}_t))$. The inflow is a boundary condition and nothing in the
  flow feeds back on it, so $g_x \equiv 0$ and $p(\bm x_t \mid \bm x_{t-1})$
  is just the AR2 prior.
\item \textbf{The mean flow is not a state.} It is a statistic of $\bm z$ with
  no equation of its own, so it is an observable, not a slow state.
\item \textbf{The fast state cannot be predicted.} Beyond a few seconds the
  turbulence at the sensors is unpredictable from $(\bm x, \bm\theta)$. A
  filter on $\bm z$ adds nothing for the slower variables. Only the
  marginalisation \eqref{eq:marginal} matters, and the ensemble spread
  approximates it.
\item \textbf{Cost.} The nested SMC--EnKF--EKF needs $N\times J$ model runs per
  slow step, and a Jacobian or $2d_z+1$ sigma points on $\bm z$. For the LES,
  neither is affordable.
\end{enumerate}

\section{Stage 1: joint ESMDA on parameters and inflow}

\subsection{Observation types}
All observation types are statistics of the same sensor time series
$u_{cs}(t_i)$, for velocity component $c$, sensor $s$ and frames
$i = 1,\dots,n$ with spacing $\Delta t$. The config chooses any non-empty subset
of them as a list, \texttt{observation.types: [bins, tke, spectral]}. The
default is \texttt{[bins]}, which reproduces current runs exactly. In the
observation vector the rows come in the order raw, bins, tke, spectral, and
within each row the sensor is the innermost index, so localisation and the
height grouping still apply. This is being implemented on branch
\texttt{feat/observation-statistics}, in
\texttt{data\_assimilation.observation\_statistics}.

\begin{center}
\begin{tabular}{@{}lll@{}}
\toprule
Type & Statistic per $(c, s)$ & Informs mainly \\
\midrule
raw & every frame $u_{cs}(t_i)$ & state, inflow \\
bins & interval means $\langle u_{cs}\rangle_k$, $k = 1,\dots,K$ & inflow $\bm x$ \\
tke & $k_{cs} = \tfrac12 \operatorname{var}_t u_{cs}$ & $\bm\theta$ (closure, $z_0$) \\
spectral & band-mean PSD $\bar E_{cs}(f_b)$, $b = 1,\dots,B$ & $\bm\theta$ (closure) \\
\bottomrule
\end{tabular}
\end{center}

\noindent \texttt{raw} and \texttt{bins} cannot be combined, because the bins
are averages of the same frames and the data would be counted twice.
\texttt{raw} with \texttt{tke} or \texttt{spectral} is allowed, but partly
redundant. The TKE of a sensor is the sum of its component entries,
$k_s = \sum_c k_{cs}$. The TKE stays split by component so that every type has
the same (component, sensor) layout.

\paragraph{Transforms.} The two turbulence statistics have optional
transforms:
\begin{align}
  \text{tke:}\quad &
  \tilde k_{cs} = \log \frac{k_{cs}}{\tfrac12 |\bar{\bm U}_s|^2}
  \quad\text{(normalisation optional, default off)}, \\
  \text{spectral:}\quad &
  \tilde E_{cs,b} = \log \frac{\bar E_{cs}(f_b)}{\operatorname{var}_t u_{cs}}
  \quad\text{(normalisation default on)},
\end{align}
where $|\bar{\bm U}_s|^2 = \sum_c \bar u_{cs}^2$ over the observed
components. The spectrum is a one-sided Welch PSD, averaged over $B$
log-spaced bands between $f_{\min}$ and the Nyquist frequency.
\begin{itemize}
\item Normalising strips the leading-order dependence on the inflow speed.
  Otherwise TKE $\propto U^2$ would push the closure parameters to absorb inflow
  errors.
\item The normalised spectrum observes only the shape of the spectrum.
\item The logarithm makes the errors closer to Gaussian, and makes a std in
  log units a relative error.
\end{itemize}

\paragraph{Update jointly, not in parallel.} The parameters affect both sets of
statistics: $z_0$ moves the street-level means, and the inflow angle moves
wakes on and off the sensors. Two separate updates would drop the
cross-covariance $C_{\bm\theta \bm h_{\mathrm{m}}}$ and its counterparts. Use one
ESMDA on the augmented vector $\bm m = (\bm\theta, \bm x)$, iterations
$a = 1,\dots,N_a$ with $\sum_a \alpha_a^{-1} = 1$:
\begin{equation}
  \bm m_j^{a+1} = \bm m_j^{a}
  + C_{\bm m \bm h}\big(C_{\bm h\bm h} + \alpha_a R\big)^{-1}
    \big(\bm y + \sqrt{\alpha_a}\,\bm\epsilon_j - \bm h(\bm m_j^a)\big),
\end{equation}
\begin{equation}
  \bm h = \begin{pmatrix}\bm h_{\mathrm{bins}}\\ \bm h_{\mathrm{tke}}\\
    \bm h_{\mathrm{spec}}\end{pmatrix},
  \quad
  R = \begin{pmatrix} R_{\mathrm{bins}} & & \\ & R_{\mathrm{tke}} & \\
    & & R_{\mathrm{spec}}\end{pmatrix}.
\end{equation}
$R$ is block-diagonal, one block per chosen type, and each block is diagonal.
This is justified because the sampling errors of a mean and of a second
moment of the same series are nearly uncorrelated. The gain then decides which
statistic informs which parameter. One diagonal block structure across
spectral bands is an approximation: neighbouring Welch bands are mildly
correlated.

\subsection{Noise model per observation type}
There are three error sources: instrument noise $\sigma_i$ (white, per raw
frame), representation error $\sigma_r$, and the sampling error of a
statistic over a finite window.

\paragraph{raw and bins (unchanged).} Each frame has variance
$\sigma_i^2 + \sigma_r^2$. For bins, the existing propagation applies:
$\sigma_i^2/m$ for a mean over $m$ frames, and $\sigma_r^2/m$ or $\sigma_r^2$ for
an independent or persistent representation error. For raw, frames
$\Delta t = 1$\,s apart are treated as independent, although the turbulence
and representation errors are correlated over seconds. This is the existing
behaviour, and it makes a raw update overconfident.

\paragraph{tke and spectral: instrument noise is a bias, added to the
prediction.} White noise adds a known amount to the expected statistics:
\begin{equation}
  \mathbb E\big[\operatorname{var}_t(u + \epsilon)\big] = \operatorname{var}_t u + \sigma_i^2,
  \qquad
  \mathbb E\big[E_{u+\epsilon}(f)\big] = E_u(f) + 2\sigma_i^2\Delta t .
\end{equation}
The observed truth series is noisy and the predictions are not, so the
bias is added to the \emph{predictions} deterministically, before normalising
and taking logs. The observed statistics are computed from the noisy series
as they are. Subtracting the bias from the observations instead would drive
the noise-dominated high-frequency bands to zero or below, where the log is
meaningless. With the bias on the prediction side, those bands look the same
in truth and prediction. They then carry no information but do no harm.

\paragraph{Sampling error.} TKE and spectra over a window of $n$ frames carry
finite-sample noise. With \texttt{std: auto} this noise is estimated from the
observed series itself, in linear units, and carried to the row's (log or
normalised) units by the delta method. It is computed on the noisy series, so
it includes the instrument part.
\begin{itemize}
\item \textbf{tke: Bartlett's formula.} With $\hat\rho(\tau)$ the sample
  autocorrelation, summed up to its first zero crossing (at most $n/4$ lags),
  \begin{equation}
    \frac{\operatorname{Var}\hat\sigma^2}{\sigma^4}
    \approx \frac{2}{n}\sum_{\tau} \hat\rho(\tau)^2 .
  \end{equation}
  With \texttt{normalize}, the error of $|\bar{\bm U}|^2$, from
  $\operatorname{Var}\bar u = (\sigma^2/n)\sum_\tau\hat\rho(\tau)$, is added in
  quadrature.
\item \textbf{spectral: spread over Welch segments.} The window PSD is the
  mean of $K$ half-overlapping segment periodograms $P_m$, so
  \begin{equation}
    \sigma_{\bar P} = \frac{\operatorname{std}_m P_m}{\sqrt K}
      \sqrt{1 + 2c\,\frac{K-1}{K}},
    \qquad c = \rho^2 \approx 0.028,
  \end{equation}
  with $\rho \approx 1/6$ the overlap correlation of adjacent Hann windows, so
  $c$ is the correlation of adjacent segment periodograms, computed from the
  window. The window must hold $K \ge 2$ segments. With
  \texttt{normalize}, the relative error of the variance is combined with the
  band's own, weighted by the band's share of the variance.
\item \textbf{Calibration.} An AR(1) Monte Carlo ($\phi = 0.9$, 120 frames)
  puts the median estimate at 0.90--1.17 of the true std for the default
  transforms. Batch means, the first approach tried, were up to 40\% off.
\item \textbf{Pooling per height.} The estimate is pooled as the median over
  sensors at the same height (\texttt{pool\_std: height}). This is refused
  unless \texttt{log} or \texttt{normalize} is on, because absolute TKE levels
  differ between sensors.
\end{itemize}

\paragraph{Instrument noise, exactly.} With constant detrending, removing the
mean also removes some noise from the lowest bins. The predicted noise PSD
is therefore $2\sigma_i^2\Delta t\,\big(1 - |W_k|^2/(N\textstyle\sum w^2)\big)$
per Welch bin $k$, with $W$ the DFT of the window: $5/6$ of the flat level at
bin 1. The observed call stores the spectral setup (frame spacing, segment
length, bands), and the predictions reuse it. A mismatch in frame count or
spacing raises an error. A constant observed series, whose logarithm would be
$-\infty$, is refused.

\paragraph{Total std.} The block std is
\begin{equation}
  \sigma = \max\!\Big(\sqrt{\sigma_{\mathrm{samp}}^2 + \sigma_{\mathrm{rep}}^2},\;
    \sigma_{\mathrm{floor}}\Big),
\end{equation}
in the block's final (log) units.
\begin{itemize}
\item $\sigma_{\mathrm{rep}}$ (\texttt{representation\_std}, default 0) is
  zero in a twin experiment. Set it for a truth from another model or from
  real sensors: a resolved LES misses the subgrid part of the TKE, for
  example.
\item A fixed \texttt{std} can replace the \texttt{auto} estimate, for example
  one calibrated from the replica noise floor of compare\_models.
\end{itemize}

\paragraph{Caveat: sampling noise is counted twice.} Each member's statistics
carry their own sampling noise $\eta_j$, which inflates $C_{\bm h\bm h}$ in the
gain. With $R$ also holding the truth's sampling variance, this noise is
counted about twice in a single update ($N_a = 1$), so the update is
conservative. That is the safe side. With $N_a > 1$ the factor is not two:
$R$ enters as $\alpha_a R$, while the members' noise stays unscaled in
$C_{\bm h\bm h}$.

\paragraph{Carry $\bm\theta$ across windows.} $\bm\theta$ is constant, so its
posterior from window $w$ is the prior for window $w+1$, while the inflow
knots $\bm x$ are re-estimated in each window. Add no random walk on
$\bm\theta$. If ensemble diversity has to be kept, use a jitter whose
variance shrinks with the ensemble size, as in the paper's
$\tilde\sigma^2 \propto N^{-3/2}$.

\paragraph{Resolution requirement.} The SGS model acts on wavelengths
$\lambda \approx 2$--$4\Delta = 4$--$8$\,m. At canopy speeds $U \approx 2$--$5$\,m/s
this means $f = U/\lambda \approx 0.3$--$1.2$\,Hz. The current output step of
1\,s gives a Nyquist frequency of 0.5\,Hz, so most of the band where
$b_0, b_1, b_2$ matter is not recorded. The sensor time series need a denser
output.

\section{Stage 2: localised ensemble smoother on the velocity field}

\subsection{No further model run}
An ensemble smoother updates trajectories that have already been simulated.
The run it needs already exists: the last forecast of each ESMDA window runs
every member with its own posterior $(\bm\theta_j, \bm x_j)$, not the posterior
mean, so the field spread includes parameter uncertainty. Stage 2 corrects those
saved frames and adds no model run. A run is needed only to forecast beyond the
window.

With the forecast ensemble $X_f \in \R^{n\times J}$ (all saved frames of all
members) and centred predicted observations $Y \in \R^{d_y\times J}$,
$Y\bm 1 = 0$, every posterior field is a combination of the member fields.

\paragraph{The gain works within $(\bm\theta, \bm x)$.} The members already
carry a posterior conditioned on $\bm y$. A cross-member gain would hold the
between-member part $\operatorname{Cov}(g(\bm\theta,\bm x)\mid\bm y)$ and use
$\bm y$ a second time on it. Let $P$ project out $[\bm 1,\ \text{parameter
anomalies}]$, in analysis space and with one row per time knot, and let
$B \in \R^{J\times\nu}$ be an orthonormal basis of its range, with
$\nu = J - 1 - p$ and $p$ the rank of the anomalies. Split
$Y = YBB\T + Y_g$, with $Y_g = Y(I - BB\T)$ the between-member part. With
$Y_w = R^{-1/2}Y$,
\begin{equation}
  X_a = X_f W,
  \quad
  W = I + B\big(\sqrt{\nu}\,M^{-1/2} - I\big)B\T
      + B M^{-1} B\T Y_w\T D_w,
  \quad
  M = \nu I + B\T Y_w\T Y_w B,
  \label{eq:es}
\end{equation}
with $D_w = R^{-1/2}D$ and $D = (\bm y - \bar{\bm y})\bm 1\T - Y_g$. This is a
deterministic transform: no observation perturbations are drawn. $D$ moves each
member to its own conditional mean, so the mean weights differ per member and
the between-member part is conditioned too. The symmetric square root gives the
within-parameter anomalies the Kalman covariance, $M^{-1}$ in ensemble
coordinates. $W \in \R^{J\times J}$ depends only on observation-space
quantities and $M$ is $\nu\times\nu$ whatever $d_y$ is. $\nu$ must stay
well above zero: a window's knots and $\bm\theta$ together count towards $p$.

\paragraph{Observations.} The observations are the raw noisy sensor frames, each
paired with the forecast frame at its time, whatever Stage 1 assimilated. Bins
are not used, and \texttt{tke} and \texttt{spectral} are excluded: they are not
linear in the frame. $R$ is diagonal and holds instrument and representation
variance only. It has no sampling error, because the target is the realised
field, not a statistic of it.

$R$ treats the frames as independent. Turbulence and representation errors
are correlated over seconds, so an update that uses several frames can be
overconfident: with AR(1) noise in the twin test of Section~\ref{sec:twin} the
spread/error ratio at the sensors fell from 0.98 to about 0.72. The run saves the lag-1
autocorrelation of the innovations per observation to show it.

\paragraph{Failed members.} A member that failed in the last forecast holds a
clone of a donor's field beside its own parameters. Such slots are excluded
from all statistics, and their columns stay unchanged.

\paragraph{What the update preserves.} The columns of $W$ sum to one, because
$B\T\bm 1 = 0$: each posterior field is an affine combination of the member
fields, with weights that can be negative. With one global $W$, every affine
constraint satisfied by all members is preserved, in particular
$\nabla\cdot\bm u = 0$ and $\bm u = 0$ inside buildings. With localisation, $W$
varies in space: a row equal across members stays exact ($\bm u = 0$ in a
building), but divergence is no longer zero. That is acceptable for
reconstruction, but a restart needs a projection first. The next window
therefore starts from the uncorrected final frame. The fields never satisfy the
nonlinear momentum equation exactly.

\subsection{Dimensionality}
\label{sec:sizes}
The sizes per case, for $u, v, w$ in float32:
\begin{center}
\begin{tabular}{@{}lrrrr@{}}
\toprule
Grid & Values per frame & Frames & Per member & $J = 32$ \\
\midrule
$50\times40\times16$ ($\Delta = 2$\,m, committed case) & $9.6\times10^4$ & 300 & 115\,MB & 3.7\,GB \\
$75\times60\times24$ ($\Delta = 1.33$\,m) & $3.2\times10^5$ & 180 & 233\,MB & 7.5\,GB \\
Barcelona $224\times224\times32$ ($\Delta = 4$\,m) & $4.8\times10^6$ & 30 & 578\,MB & 18.5\,GB \\
\bottomrule
\end{tabular}
\end{center}
The output step is 1\,s (10\,s for Barcelona). A 1\,m Xie \& Castro grid would
have $7.7\times10^5$ values per frame, but no case is configured with it. uDALES
also saves the pressure, which adds a third to the files.

The smoother \eqref{eq:es} never forms the $n\times n$ state covariance. The
state is never held in memory:
\begin{enumerate}
\item compute the weights from the observations;
\item stream the saved frames from disk one at a time, apply them, and write
  the frames back.
\end{enumerate}
Memory is $\mathcal O(J \times \text{frame})$. The saved window files of both
ESMDA save modes have the same layout, so one path serves both.

\paragraph{Localisation as a local analysis.} With $J \approx 40$, spurious
correlations dominate distant pairs, so localise in space and time. For grid
cell $\bm r$ and frame $t$, an observation at sensor $\bm r_o$ and frame time
$t_o$ enters with variance $R_o/\rho$, where
\begin{equation}
  \rho(\bm r, t; \bm r_o, t_o)
  = \operatorname{GC}(|\bm r - \bm r_o|, L_r)\,
    \operatorname{GC}(|t - t_o|, L_t),
  \quad
  W_{\bm r, t} = W\big(Y, R/\rho, \bm y\big),
  \quad
  \bm u_a(\bm r, t) = X_f(\bm r, t)\, W_{\bm r,t},
  \label{eq:local}
\end{equation}
with GC the Gaspari--Cohn taper, zero from its cutoff on. So the local
observation set is $|\bm r - \bm r_o| < L_r$ and $|t - t_o| < L_t$. The
observations are single frames, so $|t - t_o|$ is the distance between two
times. A bin would need the distance to its interval, which is not implemented.
This is a local form of the transform \eqref{eq:es}, not an LETKF: an LETKF's
anomalies keep the between-member part, and its mean weights are one vector for
all members. Each cell costs one $\nu\times\nu$ eigendecomposition. The cells
are independent and streamable, and because the transform is deterministic,
neighbouring cells cannot disagree through a sampling draw. This avoids the
in-memory gain matrix behind the $\sim$110\,GB estimate for localised
state-bearing smoothers in the case config. With no $L_r$ set, the global
\eqref{eq:es} applies, with the observations of all frames for every frame.

\paragraph{Time localisation is physically required.} The velocity decorrelates
over $\tau_{\mathrm{int}} \sim$ seconds. An observation at $t_o$ carries no
information about frames $|t - t_o| \gg \tau_{\mathrm{int}}$, so
$L_t = \mathcal O(\tau_{\mathrm{int}})$. This also keeps each local observation
set small.

\subsection{First twin test}\label{sec:twin}
One small identical-twin test, uDALES on Xie \& Castro at 2\,m ($50\times40\times16$),
$J = 32$, two 60\,s windows, one ESMDA step, $p = 6$, $\nu = 25$, four seeds
varying truth, prior, noise and inlets. There were 18 assimilated sensors and 10
held-out sensors at least 10\,m from them. The measured integral time was
$\approx 2.4$\,s. Changes are paired RMSE against the uncorrected forecast, mean
$\pm$ sd over seeds.

\begin{center}
\begin{tabular}{@{}lrrr@{}}
\toprule
Update ($L_r$\,m, $L_t$\,s) & Volume $uvw$ & Within 8\,m, $uv$ & Held-out, $uv$ \\
\midrule
Global & $+1.5 \pm 0.5\%$ & $-1.9 \pm 0.8\%$ & $-0.4 \pm 1.1\%$ \\
Local (10, 3) & $-0.9 \pm 0.1\%$ & $-10.3 \pm 0.6\%$ & $0.0\%$ \\
Local (20, 5) & $-3.1 \pm 0.2\%$ & $-18.4 \pm 0.9\%$ & $-0.5 \pm 0.3\%$ \\
Local (20, 5), $p = 0$ & $-3.3 \pm 0.2\%$ & $-19.1 \pm 0.9\%$ & $-0.5 \pm 0.3\%$ \\
Local (20, 5), shuffled $\bm y$ & $+0.5 \pm 0.4\%$ & $+4.0 \pm 1.6\%$ & $0.0 \pm 0.2\%$ \\
Local (40, 10) & $-3.0 \pm 0.6\%$ & $-17.2 \pm 0.6\%$ & $-3.7 \pm 2.9\%$ \\
\bottomrule
\end{tabular}
\end{center}

\begin{itemize}
\item Localisation is essential: the global update worsens the field away from
  the sensors, in all four seeds.
\item The local gain sits close to the sensors and has gone by 16\,m. The
  independent held-out score is much smaller than the near-sensor score.
\item $w$ barely improves, and (40, 10) worsens it by 2.5\%.
\item The projection gave no RMSE benefit in the local update here: $p = 0$ was
  slightly better in all four seeds. In the global update it limits the
  collapse of the spread (spread/error 0.44 against 0.08).
\item The local update raises the rms divergence from $10^{-7}$ to about
  $0.02$\,s$^{-1}$.
\end{itemize}

\subsection{Cheaper reduced states}
\emph{Neither option below is implemented.} Both would reuse \eqref{eq:es}, with
another state and observation operator.

\paragraph{(a) Window-mean field (not implemented).} Take the time mean of each
member's saved frames:
\begin{equation}
  \bar{\bm u}_j = \frac{1}{T}\int_{t_0}^{t_0+T} \bm u_j(t)\,dt,
  \qquad
  \bar X_f = [\bar{\bm u}_1,\dots,\bar{\bm u}_J] \in \R^{n_s\times J},
  \quad n_s = 9.6\times10^4 .
\end{equation}
The observation operator $\bar H$ interpolates the mean field to the sensors,
and the observations are the window means of the sensor series. The
update is $\bar X_a = \bar X_f \bar W$.
\begin{itemize}
\item \textbf{Cost:} $J n_s$ values, about 15\,MB for $J = 40$, so the state
  fits in memory. Spatial localisation (the form of \eqref{eq:local}) is
  affordable, and no time localisation is needed.
\item \textbf{Variant:} one mean per observation bin, $K$ fields instead of
  one, with each bin updated only from its own bin observations.
\item \textbf{Interpretation:} the mean field is roughly a deterministic
  function of the parameters plus sampling noise,
  $\bar{\bm u}_j \approx g(\bm\theta_j, \bm x_j) + \bm\varepsilon_j$. The update
  moves the mean only within the span of the member anomalies, so a structural
  model error is corrected only as far as it projects on that span, and nothing
  shows that it does. Nor can $\bm\varepsilon_j$ be assumed small: after stage 1
  the spread of $g$ is narrow, so $\bm\varepsilon_j$ may be a large share of the
  anomalies. Whether ensemble correlations are more trustworthy than for
  instantaneous fields has to be measured.
\end{itemize}

\paragraph{(b) POD coefficients (not implemented).} Expand the fluctuations of each member in a
spatial basis $\Phi_r = [\bm\phi_1,\dots,\bm\phi_r]$:
\begin{equation}
  \bm u_j(t) = \bar{\bm u} + \Phi_r\,\bm a_j(t) + \bm r_j(t),
  \qquad
  \bm a_j(t) = \Phi_r\T D\,\big(\bm u_j(t) - \bar{\bm u}\big),
\end{equation}
where $D$ is the diagonal matrix that weights the variables
(\texttt{variable\_scales}), so the modes are orthonormal in its inner product,
$\Phi_r\T D\,\Phi_r = I$. The state per member
becomes $\bm a_j(t_1{:}t_T) \in \R^{rT}$, for example $r = 100$ and
$T = 300$, giving $3\times10^4$ instead of $2.9\times10^7$ values. Per frame, the observation
operator is $H\Phi_r$ applied to the coefficients plus $H(\bar{\bm u} + \bm r_j)$.
Update each frame's coefficients from the observations within $L_t$
(time localisation is trivial in this space). Decode only the increment, as
the existing \texttt{OnlineStateReduction} does:
\begin{equation}
  \bm u_{a,j}(t) = \bm u_{f,j}(t) + \Phi_r\,\Delta\bm a_j(t),
\end{equation}
so the projection residual $\bm r_j$ is kept unchanged.

The choice of basis decides whether this adds anything:
\begin{itemize}
\item \textbf{Ensemble-anomaly basis} (the existing \texttt{svd}, rank
  $\le J-1$). This spans exactly the space of the global update, so the
  result is identical to the unlocalised smoother. It saves memory but adds
  no information.
\item \textbf{Fixed snapshot basis}, built from the window's frames
  (the \texttt{window\_snapshots} basis source) or from a long reference run. Use the method of snapshots on
  $M = JT/s$ frames with stride $s$:
  \begin{equation}
    G = S\T D S \in \R^{M\times M},\quad
    G\bm v_k = \lambda_k \bm v_k,\quad
    \bm\phi_k = S\bm v_k/\sqrt{\lambda_k}.
  \end{equation}
  Accumulate $G$ chunk by chunk from the saved frames. With $s = 10$ and $J = 32$, $M = 960$
  and $G$ is trivial in size. Truncating to the $r$ energetic modes acts as a
  \emph{scale filter}: it suppresses small-scale increments. It is not a
  substitute for spatial localisation. The energetic modes are global, so a
  spurious correlation between distant points stays in the update, and the
  sampling error of the coefficient covariance is unchanged. The reduced path
  in the code does not support localisation.
\item \textbf{Constraints are kept exactly.} Every mode is a linear
  combination of snapshots, so the modes are divergence-free and zero inside
  buildings, and so is every decoded increment. This is an advantage over
  spatial localisation.
\item \textbf{Rank is still limited by the ensemble.} The coefficient
  covariance has rank $\le J-1$ per frame, so the basis restricts where the
  update can go but does not add degrees of freedom.
\end{itemize}

\paragraph{Status.} What is implemented is the time-resolved smoother,
\eqref{eq:es} and \eqref{eq:local}. It runs on each window's posterior frames
inside the smoother, or afterwards on a finished run, and the saved posterior is
the corrected state. There are no results yet.

\section{Summary and first checks}
\begin{enumerate}
\item Joint ESMDA of $(b_0,b_1,b_2,z_0,\bm x)$ against [bin means;
  TI and normalised band spectra], with $R$ from replicas and
  $\bm\theta$ carried across windows.
\item The posterior forecast of each ESMDA window, which runs the members with
  their own posterior: no extra run.
\item Streamed, space--time-localised ensemble smoother on those frames, using
  the raw sensor frames.
\end{enumerate}
\noindent The main risk is that the SGS terms cannot be identified. Their
leverage on the sensor statistics may be small. The first checks:
\begin{itemize}
\item a paired compare\_models run varying $b_0$, to see whether the
  normalised spectra differ by more than the replica noise;
\item the posterior/prior spread ratio of each parameter after stage 1.
\end{itemize}

\end{document}
